\documentclass[11pt,onecolumn]{IEEEtran}
\usepackage{amsmath,amssymb,amsthm}
\usepackage{graphicx}
\usepackage[hidelinks]{hyperref}

\newtheorem{theorem}{Theorem}
\newtheorem{proposition}[theorem]{Proposition}
\newtheorem{corollary}[theorem]{Corollary}
\newtheorem{lemma}[theorem]{Lemma}
\newtheorem{remark}[theorem]{Remark}
\newtheorem{definition}[theorem]{Definition}

\newcommand{\E}{\mathbb{E}}
\newcommand{\R}{\mathbb{R}}
\newcommand{\tr}{\operatorname{tr}}
\newcommand{\Cov}{\operatorname{Cov}}
\newcommand{\Var}{\operatorname{Var}}
\newcommand{\hY}{\hat{\mathbf Y}}
\newcommand{\hy}{\hat{\mathbf y}}
\newcommand{\hX}{\widehat{X}}
\newcommand{\RW}{\mathcal{R}}
\newcommand{\Se}{\Sigma_{e}}
\newcommand{\Sez}{\Sigma_{e0}}
\newcommand{\Kc}{\Sigma_{T\mid S}}

\begin{document}

\title{Rate, Leakage, and Distortion for\\ Gaussian Sources with an Informed Encoder}

\author{Andrew~H.~Bond,~\IEEEmembership{Senior Member,~IEEE}%
\thanks{A. H. Bond is with the Department of Computer Engineering, San Jos\'e
State University, San Jos\'e, CA 95192 USA (e-mail: andrew.bond@sjsu.edu).
The check scripts of Appendix~\ref{app:numver} are in the author's
geometric-observation repository.}}

\markboth{Manuscript draft}{Bond: Rate, Leakage, and Distortion for Gaussian Sources with an Informed Encoder}

\maketitle

\begin{abstract}
An encoder observes a jointly Gaussian vector and describes part of it
to a decoder that sees the description alone, under a weighted quadratic
distortion. A second party holds a noisy linear observation of the
source and never decodes. One description then carries two costs: its
rate, and its conditional leakage to the second party, which is also
the least rate at which that party can re-encode it. The rate--leakage
region is known in single-letter form from secure source coding. We
determine its structure for an informed encoder, which observes more
than it must describe, when the decoder has no side information, for
every distortion weight and every observation map of the second party.
Both costs are convex in the error covariance of the description, so the
whole frontier is a convex program with unique minimizers. When the
optimal description is one-dimensional, which an explicit test
certifies, the minimum leakage is a closed-form function of the largest
generalized eigenvalue of a matrix pencil in which the distortion budget
enters, and the eigenvector is the direction the description reads. That
direction turns with the budget unless the decoder's and the second
party's geometries share it, whereas a description built for a single
party reads a fixed direction. The rate-optimal and leakage-optimal
descriptions differ whenever the second party's information does not map
the decoder's principal direction onto itself. For independent
components the leakage-optimal split of the distortion is not reverse
water-filling and can describe different components from the
rate-optimal split, and for a scalar described variable with a scalar
context the pencil gives an explicit closed form.
\end{abstract}

\begin{IEEEkeywords}
Equivocation, Gaussian sources, information leakage, rate--distortion
theory, secure source coding, side information.
\end{IEEEkeywords}

\section{Introduction}
\label{sec:intro}

One description of data is often held by more than one party. The
decoder that reconstructs from it pays its rate. A second party that
holds the description together with its own observation of the source
learns from it the conditional mutual information between the source and
the description given that observation. In a secrecy reading the second
party is an eavesdropper and this quantity is the information leakage
beyond its side information; minimizing it maximizes the eavesdropper's
equivocation about the source. In a storage reading the second party
keeps the description compressed against what it already knows, and the
same quantity is the least rate at which it can re-encode the
description. Either way one description carries two costs, set by two
different parties, and the encoder, which may observe more of the
source than either party, decides how they trade off.

The single-letter region of jointly achievable rate, distortion, and
equivocation is due to Villard and Piantanida
\cite[Theorem~3]{villard2013}, and for vector Gaussian sources Ekrem and
Ulukus \cite{ekrem2013} obtained the minimum leakage when the rate is
unconstrained, as an optimization over covariance matrices. The joint
rate--leakage region was left open: they report that they could not
solve the weighted problem that traces it \cite[p.~8]{ekrem2013}, and a
characterization through new extremal inequalities has since been
reported by Xu and Chen \cite{xuchen2021}. Without decoder side
information the region itself is elementary, because a single Gaussian
exhaustion lemma bounds both costs at once (Proposition~\ref{prop:posterior}).
What this paper contributes is the structure of that region for an
informed encoder. The
encoder observes a jointly Gaussian vector $T$, the decoder must
reproduce a linear part $Y=FT$ under a weighted quadratic distortion,
and the second party observes $S=HT+U$ for an arbitrary matrix $H$.
When $H$ reaches the parts of $T$ outside $Y$, the encoder holds
information about the second party's observation that the decoder does
not need, and this is what makes the two costs disagree.

The results are structural as well as explicit. The coding theorem
(Theorem~\ref{thm:op}) identifies the operational region with the
single-letter one, and Gaussian descriptions, indexed by their error
covariance, exhaust it (Proposition~\ref{prop:posterior}). In that
coordinate both costs are convex (Theorem~\ref{thm:convex}), so every
point of the rate--leakage frontier is the unique minimizer of a convex
program and the minimum leakage is a determinant maximization under
linear matrix inequalities. When the optimal description is
one-dimensional, which an explicit certificate decides, the minimum
leakage is
\[
L(D)=\tfrac12\log\bigl(1+\Delta/\mu_{\max}\bigr),
\]
with $\Delta$ the reduction in distortion the description must buy and
$\mu_{\max}$ the largest generalized eigenvalue of a pencil formed from
the source covariance, the decoder's weight, the conditional covariance
left after the second party's observation, and the budget
(Theorem~\ref{thm:onedim}). For a scalar described variable and a scalar
context this is an explicit closed form (Corollary~\ref{cor:scalar}).

The eigenvector of the pencil is the direction the description reads,
and it carries the paper's main structural finding. The rate-optimal
description reads the principal direction of the decoder's weighted
covariance at every budget, as do principal components, reverse
water-filling \cite{berger1971}, and the Gaussian information bottleneck
\cite{chechik2005}. The leakage-optimal description reads a direction
that turns with the budget, unless that direction is a common
generalized eigenvector of the decoder's geometry and the second party's
(Theorem~\ref{thm:turning}). No description optimized for a single party
does this (Proposition~\ref{prop:single}). The same condition governs
whether the two optimal descriptions coincide at all
(Theorem~\ref{thm:misalign}): in the one-dimensional regime they differ
whenever the second party's information does not map the decoder's
principal direction onto itself. For independent
components the leakage-optimal split of the distortion equalizes slopes
that we compute in closed form, activates components in an order that
is not the order of their variances, and can describe a disjoint set of
components from the rate-optimal split (Section~\ref{sec:parallel}).

\subsection*{Related work}

The conditional leakage studied here is the secrecy literature's
quantity. For a deterministic encoder,
$h(T^n\mid M,S^n)=h(T^n\mid S^n)-H(M\mid S^n)$, so minimizing
$H(M\mid S^n)/n$, which equals $I(T^n;M\mid S^n)/n$, is the same as
maximizing the equivocation $h(T^n\mid M,S^n)/n$. Secure lossy source
coding with side information at the eavesdropper was characterized by
Villard and Piantanida \cite{villard2013}; with no decoder side
information their Theorem~3 gives the discrete region of
Theorem~\ref{thm:dmc} below. Ekrem and Ulukus \cite{ekrem2013} treat
vector Gaussian sources with side information at both receivers under a
matrix distortion on the whole source, give an outer bound that is tight
without a rate constraint, and leave the joint region open. Our
distortion constrains a weighted part of the source and leaves the rest
unpenalized, the decoder has no side information, and the eavesdropper's
observation map is arbitrary. Xu and Chen \cite{xuchen2021} report a
characterization of the rate--distortion--equivocation function of
vector Gaussian secure source coding, and G\"unl\"u, Schaefer, Boche, and
Poor \cite{gunlu2022} evaluate the lossy rate region for scalar Gaussian
sources with a remote source and separate secrecy and privacy leakage.
Neither treats an encoder that observes components it need not
describe, a weighted partial distortion with an arbitrary observation
map, or the structure results below. Related problems price the leakage of a
description about side information
\cite{yamamoto1983,tandon2013,makhdoumi2014} or the equivocation of the
source in a cipher system \cite{yamamoto1997,schieler2014}. The leakage
endpoint is Steinberg's common-reconstruction functional
\cite{steinberg2009} on the source as the encoder observes it, an
augmentation noted for finite alphabets by Lapidoth, Mal\"ar, and Wigger
\cite{lapidoth2014}; a vector Gaussian common-reconstruction function
with every component penalized was given by Lu and Xu \cite{luxu2021}.
Weighted quadratic distortion models a decoder that cares about some
directions more than others, as in weighted rate--distortion
\cite{sakrison1968}, indirect source coding \cite{wolfziv1970}, and
task-based quantization \cite{shlezinger2019}. Parameterizing a
Gaussian description by its error covariance is a standard device in
vector Gaussian source coding \cite{rahman2012}, and the resulting
endpoint program is a determinant maximization \cite{vandenberghe1998}.

The paper is organized as follows. Section~\ref{sec:model} sets up the
model and the coding theorem, Section~\ref{sec:convex} proves the convex
structure, Section~\ref{sec:onedim} the one-dimensional regime,
Section~\ref{sec:directions} the results on directions and misalignment,
and Section~\ref{sec:parallel} independent components.
Section~\ref{sec:examples} gives instances and Section~\ref{sec:discussion}
open problems. The appendices prove the coding theorem and a Gaussian
exhaustion lemma and describe the numerical checks.

\section{Model and the Coding Theorem}
\label{sec:model}

Logarithms are base two unless written $\ln$. The encoder observes
i.i.d.\ copies of $T\in\R^p$, jointly Gaussian and centered with
$\Sigma_T\succ0$. The described variable is $Y=FT\in\R^m$ for a matrix
$F$ of full row rank, and the decoder's distortion is
$\E[(Y-\hY)^{\mathsf T}W(Y-\hY)]$ for a weight $W\succ0$. Write
\[
Q:=F^{\mathsf T}WF\succeq0,
\qquad
D_{\max}:=\tr(Q\Sigma_T).
\]
The second party, which we call the holder, observes $S=HT+U\in\R^r$
with $U\sim\mathcal N(0,\Sigma_U)$, $\Sigma_U\succ0$, independent of $T$
and of all coding randomness. Write
\[
J:=H^{\mathsf T}\Sigma_U^{-1}H\succeq0,
\qquad
\Kc=(\Sigma_T^{-1}+J)^{-1},
\]
the information that $S$ carries about $T$ and the conditional
covariance it leaves.

A code consists of an encoder $f_n:\R^{p\times n}\to\mathcal M_n$, a finite
set, and a decoder $g_n:\mathcal M_n\to\R^{m\times n}$. With $M=f_n(T^n)$
and $\hY^n=g_n(M)$, its rate, distortion, and leakage are
\[
R_n=\tfrac1n\log|\mathcal M_n|,\quad
D_n=\tfrac1n\sum_{i=1}^n\E\bigl[(Y_i-\hY_i)^{\mathsf T}W(Y_i-\hY_i)\bigr],\quad
L_n=\tfrac1nH(M\mid S^n)=\tfrac1nI(T^n;M\mid S^n).
\]
A pair $(R,L)$ is achievable at distortion $D$ when a sequence of codes
has $\limsup R_n\le R$, $\limsup L_n\le L$, and $\limsup D_n\le D$; the
closure of the achievable set is $\RW(D)$. Throughout $0<D<D_{\max}$; at
$D\ge D_{\max}$ the zero reproduction gives $R=L=0$. Write $R_{\min}(D)$
and $L(D)$ for the least rate and the least leakage over $\RW(D)$, and
call a description rate-optimal or leakage-optimal when it attains one
of them.

A test channel $p(\hy\mid t)$ is admissible when its distortion is at
most $D$; the Markov chain $\hY-T-S$ holds because the description is
formed without $S$. For a channel, let $\Sez$ and $\Se$ be the error
covariances of the best linear estimates of $T$ from $\hY$ and from
$(\hY,S)$. Appendix~\ref{app:exhaust} proves the following Gaussian
exhaustion lemma.

\begin{lemma}\label{lem:gauss}
Let $\hY\in\R^k$, for any $k$, be centered with finite second moments, let $A$ be its
regression matrix on $T$ and $N'=\hY-AT$ its residual, and suppose
$\Cov(N',S)=0$ and $\Cov N'\succ0$. Then
$I(T;\hY)\ge\tfrac12\log(\det\Sigma_T/\det\Sez)$ and
$I(T;\hY\mid S)\ge\tfrac12\log(\det\Kc/\det\Se)$, and the jointly
Gaussian channel $AT+N$, $N\sim\mathcal N(0,\Cov N')$ independent of
$(T,U)$, has the same second moments and meets both bounds with
equality.
\end{lemma}

\begin{proposition}\label{prop:posterior}
The union of the quadrants above
$\bigl(I(T;\hY),I(T;\hY\mid S)\bigr)$ over admissible channels equals the
union, over $\Sez\in\mathcal P(D)$, of the quadrants above
\begin{equation}
\bigl(r(\Sez),\ \ell(\Sez)\bigr)
=\Bigl(\tfrac12\log\frac{\det\Sigma_T}{\det\Sez},\
\tfrac12\log\frac{\det\Kc}{\det\Se}\Bigr),
\qquad
\Se^{-1}=\Sez^{-1}+J,
\label{eq:posterior}
\end{equation}
where $\mathcal P(D)=\{\Sez:\ 0\prec\Sez\preceq\Sigma_T,\ \tr(Q\Sez)\le D\}$.
\end{proposition}

\begin{proof}
Given an admissible channel, centering $\hY$ changes neither coordinate
and does not raise the distortion. If $\Cov N'$ is singular, either a
noiseless linear functional of $T$ is revealed and both coordinates are
infinite, or a nonzero combination of the components of $\hY$ vanishes
and may be dropped after an invertible linear map; so assume
$\Cov N'\succ0$. The residual is uncorrelated with $T$ by construction
and with $U$ by independence, hence with $S$, and Lemma~\ref{lem:gauss}
replaces the channel by $\hY=AT+N$ with the same distortion and no
larger coordinates. For it,
$\Sez^{-1}=\Sigma_T^{-1}+A^{\mathsf T}\Sigma_N^{-1}A$, so
$0\prec\Sez\preceq\Sigma_T$; the pair $(\hY,S)$ is a linear observation
of $T$ with independent noises, so $\Se^{-1}=\Sez^{-1}+J$; and the
distortion is at least that of $\E[Y\mid\hY]$, which is $\tr(Q\Sez)$.
Conversely, given $\Sez\in\mathcal P(D)$, factor
$\Sez^{-1}-\Sigma_T^{-1}=G^{\mathsf T}G$, let $Z=GT+N$ with
$N\sim\mathcal N(0,I)$ independent of $(T,U)$, and reproduce
$\hY=\E[Y\mid Z]$. Then $\Cov(T\mid Z)=\Sez$, the distortion is
$\tr(Q\Sez)\le D$, $\hY-T-S$ holds, and since $\hY$ is a function of $Z$,
data processing bounds its coordinates by those of $Z$, which are
\eqref{eq:posterior}.
\end{proof}

\begin{theorem}\label{thm:op}
For every $0<D<D_{\max}$, $\RW(D)$ equals the union of
Proposition~\ref{prop:posterior}.
\end{theorem}

The proof is in Appendix~\ref{app:coding}. Its discrete core is the
region of Villard and Piantanida without decoder side information
(Theorem~\ref{thm:dmc}); the Gaussian extension quantizes the side
information, the source, and the reproduction and passes to the limit.

\section{Convex Structure}
\label{sec:convex}

\begin{theorem}\label{thm:convex}
On $\{0\prec\Sez\preceq\Sigma_T\}$, the rate $r$ and the leakage $\ell$
of \eqref{eq:posterior} are strictly convex functions of $\Sez$, and the
distortion $\tr(Q\Sez)$ is linear. Consequently:
\begin{enumerate}
\item[(a)] for every $\alpha\in[0,1]$ the program
$\min\{\alpha r(\Sez)+(1-\alpha)\ell(\Sez):\Sez\in\mathcal P(D)\}$ has a
unique minimizer, the distortion constraint is active there, and the
minima over $\alpha$ trace the Pareto frontier of $\RW(D)$;
\item[(b)] in the coordinate $\Se$ the minimum leakage is a determinant
maximization,
\begin{equation}
L(D)=\tfrac12\log\det\Kc-\tfrac12\max\bigl\{\log\det\Se:\
0\prec\Se\preceq\Kc,\ \tr\bigl(Qh(\Se)\bigr)\le D\bigr\},
\qquad h(\Se):=(\Se^{-1}-J)^{-1},
\label{eq:maxdet}
\end{equation}
in which $h$ is matrix-convex and the constraint is equivalent to a
linear matrix inequality: with $J=B^{\mathsf T}B$,
$\tr(Qh(\Se))\le D$ if and only if some symmetric $Z$ has
$\tr(QZ)\le D$ and
\begin{equation}
\begin{bmatrix}
Z-\Se&\Se B^{\mathsf T}\\
B\Se&I-B\Se B^{\mathsf T}
\end{bmatrix}\succeq0 ;
\label{eq:lmi}
\end{equation}
\item[(c)] a feasible $\Se$ is the maximizer in \eqref{eq:maxdet} if and
only if $\tr(Qh(\Se))=D$ and there are $\varpi\ge0$ and $\Theta\succeq0$
with
\begin{equation}
\tfrac12\Se-\varpi\,\Sez Q\Sez=\Se\Theta\Se,
\qquad
\Theta(\Kc-\Se)=0,
\qquad \Sez=h(\Se),
\label{eq:kkt}
\end{equation}
and then $\varpi>0$.
\end{enumerate}
\end{theorem}

\begin{proof}
\emph{Convexity.} $r=\tfrac12\log\det\Sigma_T-\tfrac12\log\det\Sez$ is
strictly convex. For the leakage, write $J=B^{\mathsf T}B$ and note that
$\Se=(\Sez^{-1}+B^{\mathsf T}B)^{-1}$ is the error covariance of the best
linear estimate of a vector with covariance $\Sez$ from its noisy image
under $B$ with unit noise. For every matrix $L$, the estimator with gain
$L$ has error covariance
$E_L(\Sez)=(I-LB)\Sez(I-LB)^{\mathsf T}+LL^{\mathsf T}\succeq\Se$, with
equality at the optimal gain, and $E_L$ is affine in $\Sez$. A pointwise
minimum of affine maps in the L\"owner order, attained at each point, is
matrix-concave: at $\Sigma_t=(1-t)\Sigma_0+t\Sigma_1$ with optimal gain
$L_t$,
$\Se(\Sigma_t)=E_{L_t}(\Sigma_t)=(1-t)E_{L_t}(\Sigma_0)+tE_{L_t}(\Sigma_1)
\succeq(1-t)\Se(\Sigma_0)+t\Se(\Sigma_1)$. The map
$\Sigma\mapsto-\log\det\Sigma$ is strictly convex and decreasing in the
L\"owner order, so
$-\log\det\Se(\Sigma_t)\le-\log\det\bigl((1-t)\Se(\Sigma_0)+t\Se(\Sigma_1)\bigr)
\le(1-t)(-\log\det\Se(\Sigma_0))+t(-\log\det\Se(\Sigma_1))$, with strict
inequality in the second step when $\Sigma_0\ne\Sigma_1$, because
$\Sez\mapsto\Se$ is injective. So $\ell$ is strictly convex.

(a) $\mathcal P(D)$ is convex, the objective is strictly convex and
diverges as $\det\Sez\to0$, and the feasible set is bounded, so the
minimizer exists and is unique. Both $r$ and $\ell$ decrease when $\Sez$
increases in the L\"owner order, so at a minimizer with slack distortion
a move toward $\Sigma_T$ would stay feasible and lower the objective,
since $\Sigma_T$ itself has distortion $D_{\max}>D$. The region is
convex and upward closed, so its Pareto points are exactly the minima of
the weighted programs.

(b) The map $\Sez\mapsto\Se$ is a bijection of
$\{0\prec\Sez\preceq\Sigma_T\}$ onto $\{0\prec\Se\preceq\Kc\}$, with
inverse $h$: $\Sez\preceq\Sigma_T$ if and only if $\Se^{-1}=\Sez^{-1}+J
\succeq\Sigma_T^{-1}+J=\Kc^{-1}$. The matrix $B\Kc B^{\mathsf T}$ has the
same nonzero eigenvalues as $\Kc^{1/2}J\Kc^{1/2}\prec
\Kc^{1/2}(\Sigma_T^{-1}+J)\Kc^{1/2}=I$, so on the feasible set
$B\Se B^{\mathsf T}\prec I$, and the Woodbury identity gives
$h(\Se)=\Se+\Se B^{\mathsf T}(I-B\Se B^{\mathsf T})^{-1}B\Se$. By the
Schur complement criterion \cite[Sec.~A.5.5]{boyd2004}, $h(\Se)\preceq Z$
if and only if \eqref{eq:lmi} holds, and the block matrix there is
affine in $(\Se,Z)$; averaging two instances with $Z_i=h(\Se^{(i)})$
gives matrix convexity of $h$. Since $\tr(QZ)$ is monotone in $Z$, the
constraint $\tr(Qh(\Se))\le D$ is equivalent to the existence of $Z$ with
$\tr(QZ)\le D$ and \eqref{eq:lmi}.

(c) The program \eqref{eq:maxdet} has a strictly feasible point
($\Sez=\epsilon\Sigma_T$ with $\epsilon D_{\max}<D$), so the
Karush--Kuhn--Tucker conditions are necessary and sufficient. The
differential of $h$ is $dh=h\Se^{-1}d\Se\Se^{-1}h$, so the gradient of
$\tr(Qh)$ is $\Se^{-1}\Sez Q\Sez\Se^{-1}$, and stationarity of
$-\tfrac12\ln\det\Se+\varpi\tr(Qh)+\tr(\Theta(\Se-\Kc))$, multiplied on
both sides by $\Se$, is the first condition of \eqref{eq:kkt}, the factor
$\ln2$ being absorbed in $\varpi$. When the distortion constraint is
active its complementary slackness holds for every $\varpi\ge0$, and the
maximizer has an active constraint by (a). If $\varpi=0$, then
$\Theta=\tfrac12\Se^{-1}\succ0$ and the second condition forces
$\Se=\Kc$, which is infeasible; so $\varpi>0$.
\end{proof}

\section{The One-Dimensional Regime}
\label{sec:onedim}

Write $\Delta=D_{\max}-D$ for the reduction in distortion the description
must buy, and
\[
A(\Delta)=\Sigma_TQ\Sigma_T-\Delta\,\Sigma_T .
\]
A one-dimensional description is $Z=b^{\mathsf T}T+N$ with $b\ne0$ and
$N\sim\mathcal N(0,\sigma^2)$ independent, reproduced by $\E[Y\mid Z]$.
Its read direction is $b$, and the direction in the space of the
described variable along which it reduces the error is $F\Sigma_Tb$.

\begin{theorem}\label{thm:onedim}
(a) A one-dimensional description meets the distortion constraint with
finite leakage if and only if
$\Delta<\lambda_{\max}(W^{1/2}\Sigma_YW^{1/2})$, where
$\Sigma_Y=F\Sigma_TF^{\mathsf T}$, and the least leakage among such
descriptions is
\begin{equation}
L_{\mathrm{1d}}(D)=\tfrac12\log\Bigl(1+\frac{\Delta}{\mu_{\max}(\Delta)}\Bigr),
\label{eq:onedim}
\end{equation}
where $\mu_{\max}(\Delta)>0$ is the largest root of
$\det\bigl(A(\Delta)-\mu\Kc\bigr)=0$. It is attained with $b$ a
corresponding generalized eigenvector and
$\sigma^2=b^{\mathsf T}\Sigma_TQ\Sigma_Tb/\Delta-b^{\mathsf T}\Sigma_Tb$, at
which the distortion constraint is active.

(b) Suppose $\Delta<\lambda_{\max}(W^{1/2}\Sigma_YW^{1/2})$ and that
$\mu_{\max}(\Delta)$ is simple (if it is not, two distinct one-dimensional
optima exist, their midpoint is strictly better by
Theorem~\ref{thm:convex}, and $L(D)<L_{\mathrm{1d}}(D)$). Let $\Sez$ and
$\Se$ be the error covariances of the description of (a), and let
$\Gamma(\varpi)=\tfrac12\Se-\varpi\,\Sez Q\Sez$. Then $L(D)=L_{\mathrm{1d}}(D)$
if and only if there is $\varpi\ge0$ with $\Gamma(\varpi)\succeq0$ and
$\Gamma(\varpi)b=0$, and then
$\varpi=\tfrac12b^{\mathsf T}\Se b/(b^{\mathsf T}\Sez Q\Sez b)$.

(c) If the reproduction is scalar ($m=1$), then $L(D)=L_{\mathrm{1d}}(D)$
for every $0<D<D_{\max}$.
\end{theorem}

\begin{proof}
(a) For the description $Z$,
$\Cov(T\mid Z)=\Sigma_T-\Sigma_Tbb^{\mathsf T}\Sigma_T/(b^{\mathsf T}\Sigma_Tb+\sigma^2)$,
so the distortion of $\E[Y\mid Z]$ is
$D_{\max}-b^{\mathsf T}\Sigma_TQ\Sigma_Tb/(b^{\mathsf T}\Sigma_Tb+\sigma^2)$,
and the constraint reads
$\sigma^2\le b^{\mathsf T}\Sigma_TQ\Sigma_Tb/\Delta-b^{\mathsf T}\Sigma_Tb$.
The leakage
$I(T;Z\mid S)=\tfrac12\log\bigl(1+b^{\mathsf T}\Kc b/\sigma^2\bigr)$
decreases in $\sigma^2$, so at the largest admissible $\sigma^2$, where
the constraint is active, it equals
$\tfrac12\log\bigl(1+\Delta\,b^{\mathsf T}\Kc b/b^{\mathsf T}A(\Delta)b\bigr)$,
finite exactly when $b^{\mathsf T}A(\Delta)b>0$. Minimizing over $b$ is
maximizing the Rayleigh quotient $b^{\mathsf T}A(\Delta)b/b^{\mathsf T}\Kc b$
of a symmetric pencil with $\Kc\succ0$, whose maximum is $\mu_{\max}$.
Some $b$ has $b^{\mathsf T}A(\Delta)b>0$ if and only if
$\Sigma_T^{1/2}Q\Sigma_T^{1/2}\not\preceq\Delta I$, and the nonzero
eigenvalues of that matrix are those of $W^{1/2}\Sigma_YW^{1/2}$.

(b) The description of (a) has an active constraint and
$\Se^{-1}=\Kc^{-1}+bb^{\mathsf T}/\sigma^2$, so $\Kc-\Se$ is a positive
multiple of $\Kc bb^{\mathsf T}\Kc$ and $\Se^{-1}\Kc b$ is a positive
multiple of $b$. With $\Theta=\Se^{-1}\Gamma(\varpi)\Se^{-1}$, the
conditions \eqref{eq:kkt} at this point read $\Gamma(\varpi)\succeq0$ and
$\Gamma(\varpi)b=0$, and Theorem~\ref{thm:convex}(c) makes them necessary
and sufficient. Taking $b^{\mathsf T}\Gamma(\varpi)b=0$ gives $\varpi$.

(c) By Lemma~\ref{lem:gauss} the optimum is attained by a jointly
Gaussian channel $\hY=a^{\mathsf T}T+N$. If $\Cov(Y,\hY)=0$ its
distortion is at least $D_{\max}>D$; otherwise $\E[Y\mid\hY]$ is a
nonzero multiple of $\hY$, which leaves the leakage unchanged and does
not raise the distortion, so the optimum is a one-dimensional
description.
\end{proof}

\begin{corollary}\label{cor:scalar}
Let $T=(Y,V)$ be scalar with unit variances and correlation $\rho$,
$|\rho|<1$, let $W=1$, and let $S=V+U$ with $\Var U=\tau^2>0$; write
$s=1+\tau^2$. For $0<D<1$,
\[
L(D)=\tfrac12\log g^\star,\qquad
g^\star=\frac{(D+s-\rho^2)+\sqrt{(D+s-\rho^2)^2-4Ds(1-\rho^2)}}{2Ds},
\]
the larger root of $P(g)=Dsg^2-(D+s-\rho^2)g+(1-\rho^2)$, attained by
$\hat Y=aY+bV+N$ with $a=(g^\star-1)/g^\star$,
$b=(g^\star-1)\rho/(g^\star(g^\star s-1))$, and the noise variance that makes
the distortion equal to $D$.
\end{corollary}

\begin{proof}
By Theorem~\ref{thm:onedim}(c), $L(D)=\tfrac12\log(1+(1-D)/\mu_{\max})$
with $\Delta=1-D$. Substituting $\mu=(1-D)/(g-1)$ into
$\det(A(1-D)-\mu\Kc)$, with $\Sigma_T=\bigl[\begin{smallmatrix}1&\rho\\ \rho&1\end{smallmatrix}\bigr]$,
$Q=e_1e_1^{\mathsf T}$, and $J=e_2e_2^{\mathsf T}/\tau^2$, gives
$-(1-D)(1-\rho^2)P(g)/(s(g-1)^2)$, a direct computation confirmed by
computer algebra (Appendix~\ref{app:numver}). Since $P(1)=(D-1)\tau^2<0$
and $P(0)=1-\rho^2>0$, $P$ has one root in $(0,1)$, which corresponds to
a negative $\mu$, and one root $g^\star>1$, which corresponds to the
unique positive root $\mu_{\max}$. The read direction is the
corresponding generalized eigenvector; normalizing the reproduction
gives the stated coefficients.
\end{proof}

For a scalar described variable with vector context, Theorem
\ref{thm:onedim}(c) likewise gives $L(D)$ as an explicit function of the
largest generalized eigenvalue of a $p\times p$ pencil, $p=1+\dim V$.

\section{Directions and Misalignment}
\label{sec:directions}

\begin{proposition}\label{prop:single}
With no holder ($J=0$), the best one-dimensional description reads the
top generalized eigenvector of the pencil $(\Sigma_TQ\Sigma_T,\Sigma_T)$
at every budget, and the rate-optimal description reverse water-fills on
the eigenvalues of $W^{1/2}\Sigma_YW^{1/2}$ along fixed directions.
\end{proposition}

\begin{proof}
With $J=0$, $\Kc=\Sigma_T$ and $A(\Delta)b=\mu\Sigma_Tb$ reads
$\Sigma_TQ\Sigma_Tb=(\mu+\Delta)\Sigma_Tb$, whose eigenvectors do not
depend on $\Delta$. The rate-optimal description is the weighted
quadratic rate--distortion solution \cite{sakrison1968}: in the
coordinates $W^{1/2}Y$ the distortion is squared error, and reverse
water-filling describes fixed principal components \cite{berger1971}.
\end{proof}

The Gaussian information bottleneck has the same property: its optimal
projections are fixed eigenvectors, and its tradeoff parameter sets only
their number and scale \cite{chechik2005}. The leakage-optimal
description behaves differently.

\begin{theorem}\label{thm:turning}
Let $\mu_{\max}(\Delta)$ be simple and positive on an open interval
$\mathcal I\subseteq(0,\lambda_{\max}(W^{1/2}\Sigma_YW^{1/2}))$, with
generalized eigenvector $b(\Delta)$. Then $b(\Delta)$ spans the same line
for every $\Delta\in\mathcal I$ if and only if, for some
$\Delta_0\in\mathcal I$, $b_0=b(\Delta_0)$ satisfies
\begin{equation}
Q\Sigma_Tb_0=\lambda b_0
\quad\text{and}\quad
J\Sigma_Tb_0=\kappa b_0
\label{eq:common}
\end{equation}
for scalars $\lambda>0$ and $\kappa\ge0$, that is, if and only if $b_0$
is a common generalized eigenvector of the pencils
$(\Sigma_TQ\Sigma_T,\Sigma_T)$ and $(\Sigma_TJ\Sigma_T,\Sigma_T)$.
Otherwise $b(\Delta)$ is constant on no open subinterval of $\mathcal I$.
\end{theorem}

\begin{proof}
A simple eigenvalue of a symmetric pencil with positive definite second
matrix, and its eigenvector, are analytic in the parameter
\cite{kato1995}. Write $A(\Delta)=A_0-\Delta\Sigma_T$.

Suppose $b$ is constant on an open subinterval. Differentiating
$(A_0-\Delta\Sigma_T)b=\mu_{\max}(\Delta)\Kc b$ gives
$\Sigma_Tb=c\Kc b$ for a scalar $c$, and
$c=b^{\mathsf T}\Sigma_Tb/b^{\mathsf T}\Kc b\ge1$ because
$\Kc\preceq\Sigma_T$. Then $A_0b=(\Delta+\mu_{\max}/c)\Sigma_Tb$, that is,
$Q\Sigma_Tb=\lambda b$ with $\lambda=\Delta+\mu_{\max}/c>0$; and
multiplying $\Sigma_Tb=c\Kc b$ by $\Kc^{-1}=\Sigma_T^{-1}+J$ gives
$J\Sigma_Tb=(c-1)b$, so $\kappa=c-1\ge0$. This proves the last sentence
and the forward direction.

Conversely, if $b_0$ satisfies \eqref{eq:common}, then
$\Kc^{-1}\Sigma_Tb_0=(1+\kappa)b_0$, so $\Sigma_Tb_0=(1+\kappa)\Kc b_0$ and
$A(\Delta)b_0=(\lambda-\Delta)\Sigma_Tb_0=(\lambda-\Delta)(1+\kappa)\Kc b_0$
for every $\Delta$: $b_0$ is a generalized eigenvector with eigenvalue
$(\lambda-\Delta)(1+\kappa)$ throughout. Let
$\mathcal E=\{\Delta\in\mathcal I:(\lambda-\Delta)(1+\kappa)=\mu_{\max}(\Delta)\}$.
It contains $\Delta_0$ and is closed in $\mathcal I$ by continuity. It is
open, because near a point of $\mathcal E$ the simple eigenvalue
$\mu_{\max}$ is the only eigenvalue of the pencil close to it, and
$(\lambda-\Delta)(1+\kappa)$ is an eigenvalue close to it. Since
$\mathcal I$ is connected, $\mathcal E=\mathcal I$, and by simplicity
$b(\Delta)$ spans the line of $b_0$ throughout.
\end{proof}

\begin{corollary}\label{cor:context}
Suppose the holder observes only the context: in coordinates $T=(Y,V)$,
$F=[\,I\ 0\,]$ and $H=[\,0\ H_V\,]$. Then \eqref{eq:common} holds if
and only if $b_0=F^{\mathsf T}\beta$ with $W\Sigma_Y\beta=\lambda\beta$ and
$\Cov(S,\beta^{\mathsf T}Y)=0$.
\end{corollary}

\begin{proof}
$Q\Sigma_Tb_0=\lambda b_0$ with $\lambda>0$ puts $b_0$ in the range of
$Q$, which is the range of $F^{\mathsf T}$, so $b_0=F^{\mathsf T}\beta$
and $F^{\mathsf T}W\Sigma_Y\beta=\lambda F^{\mathsf T}\beta$. The vector
$J\Sigma_Tb_0$ lies in the range of $H^{\mathsf T}$, which meets the
range of $F^{\mathsf T}$ only at zero, so $J\Sigma_Tb_0=\kappa b_0$
forces $J\Sigma_Tb_0=0$; with $w=\Sigma_Tb_0$, $w^{\mathsf T}Jw=0$ gives
$\Sigma_U^{-1/2}Hw=0$, that is,
$\Sigma_U^{-1}H\Sigma_TF^{\mathsf T}\beta=0$, which is
$\Cov(S,\beta^{\mathsf T}Y)=0$. The converse is direct.
\end{proof}

The same condition decides whether the two optimal descriptions coincide.
Let $b_R$ be the top generalized eigenvector of
$(\Sigma_TQ\Sigma_T,\Sigma_T)$, the read direction of the best
one-dimensional rate description.

\begin{theorem}\label{thm:misalign}
Let $\lambda_1\ge\lambda_2$ be the two largest eigenvalues of
$W^{1/2}\Sigma_YW^{1/2}$ and suppose $\Delta\le\lambda_1-\lambda_2$, the
regime in which the rate-optimal description is one-dimensional,
$b_R^{\mathsf T}T+N$. If $J\Sigma_Tb_R$ is not parallel to $b_R$, the
rate-optimal description is not leakage-optimal and the leakage-optimal
description is not rate-optimal, so each pays a strictly positive excess
in the other cost. If $J\Sigma_Tb_R$ is parallel to $b_R$, then $b_R$ is a
generalized eigenvector of the leakage pencil at every $\Delta$, and
inside the one-dimensional regime of Theorem~\ref{thm:onedim} the two
descriptions coincide exactly when it is the top one.
\end{theorem}

\begin{proof}
In the stated regime reverse water-filling describes only the principal
component, so the rate-optimal description is the one-dimensional
description reading $b_R$, where $Q\Sigma_Tb_R=\lambda_1b_R$. If it were
leakage-optimal, then $L(D)=L_{\mathrm{1d}}(D)$ and $b_R$ would maximize
the Rayleigh quotient of Theorem~\ref{thm:onedim}(a), hence be a
generalized eigenvector of $(A(\Delta),\Kc)$, that is,
$(\lambda_1-\Delta)\Sigma_Tb_R\propto\Kc b_R$; multiplying by
$\Kc^{-1}=\Sigma_T^{-1}+J$ gives $J\Sigma_Tb_R\propto b_R$. So when
$J\Sigma_Tb_R$ is not parallel to $b_R$ the rate-optimal description is not
leakage-optimal, and since both minimizers of Theorem~\ref{thm:convex}(a)
are unique, the leakage-optimal description is not rate-optimal either.
The same computation read backward shows that $J\Sigma_Tb_R\propto b_R$
makes $b_R$ a generalized eigenvector at every $\Delta$, and inside the
one-dimensional regime the leakage-optimal description reads the top
one.
\end{proof}

In the scalar case of Corollary~\ref{cor:scalar}, $b_R=e_1$ and
$J\Sigma_Te_1=(\rho/\tau^2)e_2$, which is parallel to $e_1$ only when
$\rho=0$. So the two descriptions differ exactly when the context is
correlated with the described variable. The rate-optimal channel is the
classical reverse channel, with leakage
$\tfrac12\log\bigl[((1-D)(1-\rho^2/s)+D)/D\bigr]$, strictly above $L(D)$,
and the leakage-optimal channel of Corollary~\ref{cor:scalar} has rate
strictly above $\tfrac12\log(1/D)$. The rate excess is unbounded toward the
clean-context boundary $\tau^2\to0$ at $D=1-\rho^2$, where the
leakage-optimal channel pins the reproduction to the context.

\begin{corollary}\label{cor:nonred}
There are sources for which the leakage-optimal one-dimensional read
takes different directions at two budgets. For such a source no
description problem for a single party with a fixed weighted quadratic
distortion has the leakage-optimal descriptions as its optimal
descriptions at both budgets.
\end{corollary}

\begin{proof}
Take the scalar source of Corollary~\ref{cor:scalar} with $\rho\ne0$. By
Theorem~\ref{thm:onedim}(c) the leakage-optimal description is the
one-dimensional description at every $D$, and its read direction is
proportional to $(g^\star s-1,\rho)$, the ratio of the coefficients in
Corollary~\ref{cor:scalar}. Since $\ell$ is strictly decreasing
(Lemma~\ref{lem:slope}), $g^\star$ is strictly decreasing in $D$, so the
read line differs at any two budgets. By Proposition~\ref{prop:single},
applied to any fixed weight on $T$, the best one-dimensional description
for a single party reads the same line at both budgets. The coupled
instance of Section~\ref{sec:examples}, whose leakage-optimal
descriptions reduce the error along $47.48^\circ$ and $41.15^\circ$ at
two budgets, illustrates the same effect in the space of the described
variable.
\end{proof}

\section{Independent Components}
\label{sec:parallel}

Suppose $Y=(Y_1,\dots,Y_m)$, $W$ is diagonal with entries $w_j$, and each
$Y_j$ has its own context $V_j$ and holder observation $S_j=V_j+U_j$, the
triples $(Y_j,V_j,S_j)$ being independent, with $\Var Y_j=\sigma_j^2$,
$\Var V_j=1$, $\E[Y_jV_j]=\rho_j\sigma_j$, $|\rho_j|<1$, and
$\Var U_j=\tau_j^2>0$. Absorb $w_j$ into $\sigma_j^2$. Write $\ell_j$ for
the scalar function $\delta\mapsto\tfrac12\log g^\star(\delta)$ of
Corollary~\ref{cor:scalar} at $(\rho_j,\tau_j^2)$, extended by zero for
$\delta\ge1$, and $L_j(D_j)=\ell_j(D_j/\sigma_j^2)$.

\begin{lemma}\label{lem:slope}
For $0<\delta<1$, with $g=g^\star(\delta)$, $s=1+\tau^2$, $k=gs-1$, and
$\Delta_\delta=\sqrt{(\delta+s-\rho^2)^2-4\delta s(1-\rho^2)}>0$,
\[
\ell'(\delta)=-\frac{k}{2\ln2\,\Delta_\delta},
\qquad
\ell''(\delta)=\frac{(1-\rho^2)k^3
+\rho^2\tau^2g^2\bigl(\tau^4g+\tau^2(2g+1)+g+1\bigr)}{2\ln2\;g\,\Delta_\delta^3}>0,
\]
so $\ell$ is strictly convex and decreasing, $\ell'\to-\infty$ as
$\delta\to0$, and $\ell'\to-\tau^2/(2\ln2(\rho^2+\tau^2))$ as $\delta\to1$.
\end{lemma}

\begin{proof}
The larger root is simple, with $\partial_gP=\Delta_\delta$ there, so
$g^\star$ is smooth; differentiating $P(g^\star(\delta))=0$ with
$\partial_\delta P=gk$ gives $g'=-gk/\Delta_\delta$ and the first formula.
Differentiating once more and eliminating $\delta$ through
$\delta gk=(s-\rho^2)g-(1-\rho^2)$ gives the second, an identity of
polynomials confirmed by computer algebra (Appendix~\ref{app:numver});
its numerator is positive. The limits follow from $g^\star\to\infty$,
$\Delta_\delta\to s-\rho^2$ as $\delta\to0$ and $g^\star\to1$, $k\to\tau^2$,
$\Delta_\delta\to\rho^2+\tau^2$ as $\delta\to1$.
\end{proof}

Define the activation level
$\kappa_j=-L_j'(\sigma_j^2-)=\tau_j^2/(2\ln2\,\sigma_j^2(\rho_j^2+\tau_j^2))$;
for the rate it would be $1/(2\ln2\,\sigma_j^2)$.

\begin{theorem}\label{thm:parallel}
(a) $\RW(D)$ is the union, over $D_j\ge0$ with $\sum_jD_j\le D$, of the
Minkowski sums of the scalar regions of the components, attained by
product channels.
(b) $L(D)=\min\sum_jL_j(D_j)$ over such allocations, with a unique
minimizer $D_j=D_j(\theta^\star)$, where $D_j(\theta)=\sigma_j^2$ for
$\theta\le\kappa_j$ and otherwise solves $-L_j'(D_j)=\theta$, and
$\theta^\star$ is the unique level with $\sum_jD_j(\theta^\star)=D$.
(c) As $D$ decreases, the leakage-optimal description begins to describe
component $j$ in decreasing order of $\sigma_j^2(1+\rho_j^2/\tau_j^2)$,
while the rate-optimal description proceeds in decreasing order of
$\sigma_j^2$.
\end{theorem}

\begin{proof}
(a) By Theorem~\ref{thm:op} it suffices to identify the single-letter
union. For an admissible channel with $D_j=\E(Y_j-\hat Y_j)^2$, each
induced channel $p(\hat y_j\mid t_j)$ satisfies $\hat Y_j-T_j-S_j$, and
since the pairs $(T_j,S_j)$ are independent the chain rule gives
$I(T;\hY)\ge\sum_jI(T_j;\hat Y_j)$ and
$I(T;\hY\mid S)=\sum_j[I(T_j;\hY,S,T^{j-1})-I(T_j;S_j)]\ge\sum_jI(T_j;\hat Y_j\mid S_j)$.
A product of component channels attains the sum with equality.
(b) Each $L_j$ is convex and nonincreasing, strictly convex on
$(0,\sigma_j^2)$ by Lemma~\ref{lem:slope}, with left derivative
$-\kappa_j$ at $\sigma_j^2$. An allocation is optimal if and only if some
$\theta\ge0$ lies in $-\partial L_j(D_j)$ for every $j$ with
$\theta(\sum_jD_j-D)=0$. Since $D<\sum_j\sigma_j^2$, $\theta>0$; no $D_j$
is $0$ because $-L_j'\to\infty$ there; and $\sum_jD_j(\theta)$ decreases
continuously and strictly from $\sum_j\sigma_j^2$ at $\min_j\kappa_j$ to
$0$, so $\theta^\star$ is unique. (c) $D_j(\theta)<\sigma_j^2$ exactly when
$\theta>\kappa_j$, and $\theta^\star$ increases as $D$ decreases.
\end{proof}

\begin{corollary}\label{cor:disjoint}
If $m=2$, $\sigma_1^2>\sigma_2^2$, and
$\sigma_1^2(1+\rho_1^2/\tau_1^2)<\sigma_2^2(1+\rho_2^2/\tau_2^2)$, then for every
$D\in[\max\{2\sigma_2^2,\sigma_1^2+D_2(\kappa_1)\},\sigma_1^2+\sigma_2^2)$ the
rate-optimal description describes component $1$ alone and the
leakage-optimal description describes component $2$ alone.
\end{corollary}

\begin{proof}
Reverse water-filling describes component $1$ alone exactly for
$2\sigma_2^2\le D<\sigma_1^2+\sigma_2^2$; Theorem~\ref{thm:parallel}(b)
describes component $2$ alone exactly when $\kappa_2<\theta^\star\le\kappa_1$,
that is, for $\sigma_1^2+D_2(\kappa_1)\le D<\sigma_1^2+\sigma_2^2$.
\end{proof}

\section{Instances}
\label{sec:examples}

\emph{Scalar.} At $(\rho^2,\tau^2,D)=(0.75,0.5,0.3)$ the rate-optimal
channel pays $0.0349$ bits of excess leakage and the leakage-optimal
channel $0.0400$ bits of excess rate; over the box
$\rho^2\in[0.05,0.95]$, $\tau^2\in[\tfrac14,4]$, $D\in[0.05,0.95]$ the
rate excess of the leakage-optimal channel stays below $0.114$ bits and
the leakage excess of the rate-optimal channel below $0.077$ bits, and at
$(\rho^2,\tau^2,D)=(0.5,10^{-3},0.5)$ the rate excess reaches $1.537$ bits.
Two sources with the same law of $(Y,S)$ up to scaling of $S$,
$(\rho^2,\tau^2)=(\tfrac35,\tfrac15)$ and $(\tfrac34,\tfrac12)$, have
$L=1.1880$ and $1.2105$ bits at $D=0.1$: the leakage is not determined by
the law of $(Y,S)$, because the encoder observes the context.

\emph{Coupled.} Take $\Sigma_Y=\mathrm{diag}(1,0.4)$, $W=I$, a scalar
context with $\Var V=1$ and $\Cov(Y,V)=0.6\,u$ for the unit vector $u$ at
$50^\circ$, and $S=V+U$ with $\Var U=0.05$. The certificate of
Theorem~\ref{thm:onedim}(b) holds for $D\ge1.08530$, out of
$D_{\max}=1.4$. The rate-optimal description reduces the error along the
first axis for every $D\ge0.8$. The leakage-optimal description reduces
it along $47.48^\circ$, $46.76^\circ$, $45.07^\circ$, and $41.15^\circ$ at
$D=1.38$, $1.3$, $1.2$, and $1.1$. Four sources with the same $u$ and the same law of
$(Y,S)$ up to scaling, $(|\Cov(Y,V)|^2,\tau^2)=(0.33,0.1)$, $(0.36,0.2)$,
$(0.39,0.3)$, $(0.42,0.4)$, read along $42.83^\circ$, $37.61^\circ$,
$33.75^\circ$, and $30.84^\circ$ at $D=1.3$, each inside its certificate
regime: the direction, too, is not
determined by the law of $(Y,S)$.

\emph{Independent components.} With
$(\sigma_1^2,\rho_1^2,\tau_1^2)=(1,0,1)$ and
$(\sigma_2^2,\rho_2^2,\tau_2^2)=(0.6,0.9,0.05)$,
$\kappa_1=0.7213$, $\kappa_2=0.0633$, and the two descriptions describe
disjoint components for every $D\in[1.2002,1.6)$. At $D=1.4$ the
rate-optimal description pays $R=L=0.1610$ bits on component $1$ and the
leakage-optimal description pays $L=0.0197$ and $R=0.3272$ bits on
component $2$.

\section{Discussion}
\label{sec:discussion}

A description held by two parties has a cost for each, and the
encoder's observations decide how the costs trade off. For Gaussian
sources without decoder side information the region is now explicit in
structure: a convex program in the error covariance, a closed form in a
one-dimensional regime, and directions that turn with the budget exactly
when the two parties' geometries are not aligned. Several problems remain
open. With side information at the decoder the weighted problem of
Ekrem and Ulukus was left open there, and a solution is reported by Xu
and Chen \cite{xuchen2021}; it is harder than the present case because the Gaussian exhaustion lemma
no longer handles both costs at once once the decoder bins against its
own observation. A closed form below the one-dimensional regime, and for
its boundary as a function of the source, is not known. Stationary
sources and finite blocklengths are not treated.

\appendices

\section{The Coding Theorem}
\label{app:coding}

\subsection{Finite alphabets}

Let $(X,S)$ be a discrete memoryless pair with finite alphabets and let
$d$ be a distortion with values in $[0,d_{\max}]$; codes, rate,
distortion, and leakage $H(M\mid S^n)/n$ are as in Section~\ref{sec:model}.

\begin{theorem}[{\cite[Theorem~3]{villard2013}, without decoder side information}]
\label{thm:dmc}
The closure of the achievable set at distortion $D$ is
$\bigcup_{p(\hat x\mid x):\,\E d(X,\hX)\le D}\{(R,L):R\ge I(X;\hX),\ L\ge I(X;\hX\mid S)\}$,
and the union is convex and closed.
\end{theorem}

\begin{proof}[Proof outline]
Villard and Piantanida prove the region in the equivocation form
$\Delta\le H(X\mid V)+0-I(X;S\mid U)$, $U-V-X-S$, with no decoder side
information; taking $U=V$ is optimal, because
$I(X;S\mid U)=I(V;S\mid U)+I(X;S\mid V)\ge I(X;S\mid V)$, and
$\Delta=H(X\mid S)-L$ for a deterministic encoder, which gives the
displayed form with $V$ in place of $\hX$, and a reproduction that is a
function of $V$ only lowers both coordinates. For the Gaussian extension
we use the following finite-blocklength form of the achievability,
obtained by covering at rate $I(X;\hX)+2\epsilon$, a random binning of the
codeword index into $2^{n(I(X;\hX\mid S)+3\epsilon)}$ bins used only as a
proof device, Fano's inequality for a hypothetical decoder holding the
bin index and $S^n$, and selection of one deterministic code: for every
channel with finite reproduction alphabet and every $\epsilon>0$ there
are $\epsilon_n\to0$ and, for all large $n$, a deterministic code with
$|\mathcal M_n|\le2^{n(I(X;\hX)+2\epsilon)}$, distortion at most
$\E d(X,\hX)+\epsilon+2\epsilon_nd_{\max}$, and
$H(M\mid S^n)\le n(I(X;\hX\mid S)+3\epsilon)+1+2\epsilon_nn(I(X;\hX)+2\epsilon)$.
Convexity of both coordinates in the channel, for the conditional one by
the decomposition $I(X;\hX\mid S)=\sum_sp(s)I_{p(x\mid s)}(q)$, makes the
union convex without time sharing.
\end{proof}

\subsection{The Gaussian source}

\begin{proof}[Proof of Theorem~\ref{thm:op}]
\emph{Converse.} For any code, $\log|\mathcal M_n|\ge I(T^n;\hY^n)\ge
\sum_iI(T_i;\hY_i)$ and $H(M\mid S^n)\ge I(T^n;\hY^n\mid S^n)
=\sum_i[I(T_i;\hY^n,T^{i-1},S^n)-I(T_i;S_i)]\ge\sum_iI(T_i;\hY_i\mid S_i)$,
by the chain rule and the independence of the pairs $(T_i,S_i)$; every
term is finite, being at most $I(T_i;S_i)+\log|\mathcal M_n|$ with
$I(T_i;S_i)=\tfrac12\log\det(H\Sigma_TH^{\mathsf T}+\Sigma_U)/\det\Sigma_U$.
The mixture of the per-letter channels satisfies the Markov chain, has
distortion $D_n$, and by convexity of both coordinates in the channel
(for general alphabets the conditional coordinate is an integral over
$s$ of relative entropies affine in the channel \cite{csiszar1967,wyner1978})
lies below $(R_n,L_n)$, so $(R_n,L_n)$ is in the union at $D_n$.

\emph{Closedness.} Let $(R_j,L_j)\to(R,L)$ lie in the union at levels
$D_j$ with $\limsup D_j\le D$, with $\Sigma_j\in\mathcal P(D_j)$ below them
by Proposition~\ref{prop:posterior}. The $\Sigma_j$ lie in the compact
set $0\preceq\Sigma\preceq\Sigma_T$ and $\det\Sigma_j\ge\det\Sigma_T\,2^{-2R_j}$
is bounded away from zero, so a subsequence converges to some
$\Sigma^\star\succ0$ in $\mathcal P(D)$; $r$ and $\ell$ are continuous on
positive definite matrices, so $(R,L)$ lies above
$(r(\Sigma^\star),\ell(\Sigma^\star))$.

\emph{Achievability.} Fix an admissible channel with $I(T;\hY)<\infty$;
$\E\|\hY\|^2<\infty$ because its distortion is finite and $W\succ0$. Let $q_k$ be the saturating dyadic quantizer of
level $k$ applied coordinatewise: the bounded cells
$[j2^{-k},(j+1)2^{-k})$, $-4^k\le j<4^k$, map to their left endpoints and
the two tail cells to $\pm2^k$; the partitions are nested and generate
the Borel $\sigma$-field.

\emph{Step A.} The $\sigma$-fields of $S_\eta:=q_\eta(S)$ increase to
$\sigma(S)$, so $I(\hY;S_\eta)\uparrow I(\hY;S)$ \cite[Ch.~2]{pinsker1964};
fix $\eta$ with $I(\hY;S_\eta)\ge I(\hY;S)-\epsilon$. Since $\hY-T-S_\eta$,
$I(T;\hY\mid S_\eta)=I(T;\hY)-I(\hY;S_\eta)\le I(T;\hY\mid S)+\epsilon$.

\emph{Step B.} With $T_k=q_k(T)$ and $\hY_k=q_k(\hY)$, the surrogate
source $(T_k,S_\eta)$ is discrete memoryless, the surrogate channel is
$p(\hy\mid t_k)=\Pr(\hY_k=\hy\mid T_k=t_k)$, and the surrogate distortion
$d_k(t_k,\hy)=\E[(Y-\hy)^{\mathsf T}W(Y-\hy)\mid T_k=t_k]$ is bounded for
each $k$, every cell having positive probability since
$\Sigma_T\succ0$. The finite-blocklength achievability of
Theorem~\ref{thm:dmc} gives a code of rate at most
$I(T_k;\hY_k)+2\epsilon\le I(T;\hY)+2\epsilon$. Its true distortion equals
its surrogate distortion, because given $T_k^n$ the law of $Y_i$ depends
on $T_{k,i}$ alone. The surrogate distortion of the channel converges to
$\E[(Y-\hY)^{\mathsf T}W(Y-\hY)]\le D$, because $\E[Y\mid T_k]\to Y$ in
$L^2$ by martingale convergence and $\hY_k\to\hY$ in $L^2$. The surrogate
law of $(\hY_k,S_\eta)$ converges weakly to the true law of
$(\hY,S_\eta)$, by the bounded-convergence and martingale argument
applied to test functions bounded on $\R^m\times q_\eta(\R^r)$ and
continuous in the first argument, and the Markov chain $\hY-T-S$. Mutual
information is lower semicontinuous under weak convergence
\cite{posner1975}, so
$I_k(T_k;\hY_k\mid S_\eta)=I(T_k;\hY_k)-I_k(\hY_k;S_\eta)\le I(T;\hY\mid S)+2\epsilon$
for large $k$, and the code has $H(M\mid S_\eta^n)\le n(I(T;\hY\mid S)+5\epsilon+o(1))$.

\emph{Step C.} $S_\eta^n$ is a function of $S^n$, so
$H(M\mid S^n)\le H(M\mid S_\eta^n)$. A diagonal sequence
$\epsilon\downarrow0$ achieves the corner
$(I(T;\hY),I(T;\hY\mid S))$ at distortion $D$, and with the converse,
closedness, and Proposition~\ref{prop:posterior} the region is as stated.
\end{proof}

\section{The Gaussian Exhaustion Lemma}
\label{app:exhaust}

\begin{proof}[Proof of Lemma~\ref{lem:gauss}]
\emph{The divergence bound.} Conditional mutual information is the
expected divergence of the conditional law of $T$ given $(\hY,S)$ from
the conditional law given $S$ \cite{wyner1978,pinsker1964}. Fix a
realization, let $P$ be the conditional law with mean $\mu$ and
covariance $C$, $P_G=\mathcal N(\mu,C)$, and
$Q_S=\mathcal N(\mu_S,\Kc)$. If $P$ has no density the divergence is
infinite. Otherwise $\log(dP_G/dQ_S)$ is quadratic in $t$ and $P$, $P_G$
share two moments, so $D(P\|Q_S)=D(P\|P_G)+D(P_G\|Q_S)\ge D(P_G\|Q_S)$.
Averaging $D(P_G\|Q_S)$, the trace term plus the mean term,
$\E[\tr(\Kc^{-1}C)+(\mu-\mu_S)^{\mathsf T}\Kc^{-1}(\mu-\mu_S)]$, averages to
$p$ by the law of total covariance, and by concavity of $\ln\det$ and Jensen,
$\E[-\ln\det C]\ge-\ln\det\E C\ge-\ln\det\Se$, since the conditional mean
has error covariance no larger than any linear estimator. Hence
$I(T;\hY\mid S)\ge\tfrac12\log(\det\Kc/\det\Se)$, and the same argument
without $S$ gives $I(T;\hY)\ge\tfrac12\log(\det\Sigma_T/\det\Sez)$.

\emph{The determinant identity.} With $\hY=AT+N'$, $\Cov(N',T)=0$, and
$\Cov(N',S)=0$, the change of variables $(T,S,\hY)\mapsto(T,S,N')$ is
block triangular with unit diagonal, so
$\det\Cov(T,S,\hY)=\det\Sigma_S\det\Kc\det\Sigma_N$; the Schur factorization
in the other order gives $\det\Cov(\hY,S)\det\Se$, and
$\Cov(\hY)-\Cov(\hY,S)\Sigma_S^{-1}\Cov(S,\hY)=A\Kc A^{\mathsf T}+\Sigma_N\succ0$.
So $\det\Se=\det\Kc\det\Sigma_N/\det(A\Kc A^{\mathsf T}+\Sigma_N)$, finite and
positive.

\emph{Equality.} For the jointly Gaussian channel $AT+N$ the conditional
law of $T$ given $(\hY,S)$ is Gaussian with the linear estimate as mean
and $\Se$ as covariance, so every inequality is an equality; its second
moments match the given channel's.
\end{proof}

\section{Numerical Checks}
\label{app:numver}

The identities and instances were checked independently of the
derivations by scripts run on a remote workstation.
\texttt{verify\_general.py} checks, for random sources with general
weights and observation maps, that the pencil value of
Theorem~\ref{thm:onedim} equals the brute-force optimum over Gaussian
descriptions whenever that optimum is one-dimensional, that both costs
are convex in the error covariance (16000 midpoint tests), that the
weighted minimizers are unique, the two directions of
Theorem~\ref{thm:turning}, and the criterion of
Theorem~\ref{thm:misalign}. \texttt{verify\_coupled.py},
\texttt{verify\_parallel.py}, and \texttt{verify\_pencil.py} check the
scalar identity of Corollary~\ref{cor:scalar} symbolically, the
determinant program against brute force, the certificate, Lemma
\ref{lem:slope}, the allocation, and the numbers of
Section~\ref{sec:examples}. No proof depends on them; the instances
illustrate theorems proved in the text.

\section*{Acknowledgment}

In accordance with the IEEE policy on AI-generated text, the author
discloses the use of Anthropic's Claude models (Claude Opus 5.5) through
the Claude Code environment in 2026, for drafting prose and proofs from
the author's statements of the results, for the prior-art search, for
the check scripts, and for independent referee passes on the proofs. The
author is responsible for the content.

\begin{thebibliography}{99}
\bibitem{villard2013} J. Villard and P. Piantanida, ``Secure multiterminal
source coding with side information at the eavesdropper,'' \emph{IEEE
Trans. Inf. Theory}, vol. 59, no. 6, pp. 3668--3692, 2013.
\bibitem{ekrem2013} E. Ekrem and S. Ulukus, ``Secure lossy transmission of
vector Gaussian sources,'' \emph{IEEE Trans. Inf. Theory}, vol. 59, no. 9,
pp. 5466--5487, 2013.
\bibitem{berger1971} T. Berger, \emph{Rate Distortion Theory: A Mathematical
Basis for Data Compression}. Englewood Cliffs, NJ: Prentice-Hall, 1971.
\bibitem{chechik2005} G. Chechik, A. Globerson, N. Tishby, and Y. Weiss,
``Information bottleneck for Gaussian variables,'' \emph{J. Mach. Learn.
Res.}, vol. 6, pp. 165--188, 2005.
\bibitem{xuchen2021} Y. Xu and G. Chen, ``New proofs of extremal
inequalities with applications,'' arXiv:2109.00377, 2021.
\bibitem{gunlu2022} O. G\"unl\"u, R. F. Schaefer, H. Boche, and H. V. Poor,
``Secure and private source coding with private key and decoder side
information,'' arXiv:2205.05068, 2022; shorter version in \emph{Proc. IEEE
Inf. Theory Workshop (ITW)}, 2022. % verify journal version (Entropy) on final
\bibitem{yamamoto1983} H. Yamamoto, ``A source coding problem for sources
with additional outputs to keep secret from the receiver or wiretappers
(Corresp.),'' \emph{IEEE Trans. Inf. Theory}, vol. 29, no. 6, pp. 918--923,
1983.
\bibitem{tandon2013} R. Tandon, L. Sankar, and H. V. Poor,
``Discriminatory lossy source coding: Side information privacy,''
\emph{IEEE Trans. Inf. Theory}, vol. 59, no. 9, pp. 5665--5677, 2013.
\bibitem{makhdoumi2014} A. Makhdoumi, S. Salamatian, N. Fawaz, and
M. M\'edard, ``From the information bottleneck to the privacy funnel,''
in \emph{Proc. IEEE Inf. Theory Workshop (ITW)}, 2014, pp. 501--505.
\bibitem{yamamoto1997} H. Yamamoto, ``Rate-distortion theory for the Shannon
cipher system,'' \emph{IEEE Trans. Inf. Theory}, vol. 43, no. 3, pp. 827--835,
1997.
\bibitem{schieler2014} C. Schieler and P. Cuff, ``Rate-distortion theory for
secrecy systems,'' \emph{IEEE Trans. Inf. Theory}, vol. 60, no. 12,
pp. 7584--7605, 2014.
\bibitem{steinberg2009} Y. Steinberg, ``Coding and common reconstruction,''
\emph{IEEE Trans. Inf. Theory}, vol. 55, no. 11, pp. 4995--5010, 2009.
\bibitem{lapidoth2014} A. Lapidoth, A. Mal\"ar, and M. Wigger, ``Constrained
source-coding with side information,'' \emph{IEEE Trans. Inf. Theory},
vol. 60, no. 6, pp. 3218--3237, 2014.
\bibitem{luxu2021} J. Lu and Y. Xu, ``Vector Gaussian Heegard-Berger problem
with common reconstructions,'' in \emph{Proc. Int. Conf. Wireless Commun.
Signal Process. (WCSP)}, 2021, pp. 1--5.
\bibitem{sakrison1968} D. J. Sakrison, ``The rate distortion function of a
Gaussian process with a weighted square error criterion,'' \emph{IEEE Trans.
Inf. Theory}, vol. 14, no. 3, pp. 506--508, 1968.
\bibitem{wolfziv1970} J. K. Wolf and J. Ziv, ``Transmission of noisy
information to a noisy receiver with minimum distortion,'' \emph{IEEE
Trans. Inf. Theory}, vol. 16, no. 4, pp. 406--411, 1970.
\bibitem{shlezinger2019} N. Shlezinger, Y. C. Eldar, and M. R. D. Rodrigues,
``Hardware-limited task-based quantization,'' \emph{IEEE Trans. Signal
Process.}, vol. 67, no. 20, pp. 5223--5238, 2019.
\bibitem{rahman2012} M. S. Rahman and A. B. Wagner, ``Rate region of the
Gaussian scalar-help-vector source-coding problem,'' \emph{IEEE Trans. Inf.
Theory}, vol. 58, no. 1, pp. 172--188, 2012.
\bibitem{vandenberghe1998} L. Vandenberghe, S. Boyd, and S.-P. Wu,
``Determinant maximization with linear matrix inequality constraints,''
\emph{SIAM J. Matrix Anal. Appl.}, vol. 19, no. 2, pp. 499--533, 1998.
\bibitem{boyd2004} S. Boyd and L. Vandenberghe, \emph{Convex Optimization}.
Cambridge, U.K.: Cambridge Univ. Press, 2004.
\bibitem{kato1995} T. Kato, \emph{Perturbation Theory for Linear Operators},
reprint of the 1980 ed. Berlin, Germany: Springer, 1995.
\bibitem{pinsker1964} M. S. Pinsker, \emph{Information and Information
Stability of Random Variables and Processes}. San Francisco, CA:
Holden-Day, 1964.
\bibitem{posner1975} E. C. Posner, ``Random coding strategies for minimum
entropy,'' \emph{IEEE Trans. Inf. Theory}, vol. 21, no. 4, pp. 388--391,
1975.
\bibitem{wyner1978} A. D. Wyner, ``A definition of conditional mutual
information for arbitrary ensembles,'' \emph{Inf. Control}, vol. 38, no. 1,
pp. 51--59, 1978.
\bibitem{csiszar1967} I. Csisz\'ar, ``Information-type measures of difference
of probability distributions and indirect observations,'' \emph{Studia Sci.
Math. Hungar.}, vol. 2, pp. 299--318, 1967.
\end{thebibliography}

\end{document}
